%%%%%%
%
% $Author: Yogesh,Rakesh,Abhishek $
% $Date: 29.04.2024 $
% $Path: BA24-02-Time-Series\report\Contents\en$
% $Version: 1.0 $
% $Reviewed by: Yogesh,Rakesh,Abhishek $
%
% !TeX encoding = utf8
% !TeX root = Rename
% !TeX TXS-program:bibliography = txs:///bibtex
%
%%%%%%

\chapter{Data Description}


The dataset used in this project contains comprehensive air quality data spanning from the year 2010 to 2023 for 453 cities across India. This dataset provides a comprehensive overview of the air quality conditions in various Indian cities, covering a wide geographical distribution and a substantial time period. The data has been sourced from the Central Control Room for Air Quality Management, making it a reliable and authoritative source of information \cite{CPCB:2023}.

\section{Dataset Characteristics}

\subsection{Time Period}
The dataset covers a span of 13 years, from 2010 to 2023, allowing for the analysis of long-term air quality trends.

\subsection{Geographic Scope}
It encompasses 453 cities, representing various regions of India and providing a diverse set of air quality conditions.

\subsection{Parameters}
The dataset includes key air quality parameters such as PM2.5 (particulate matter 2.5 micrometers or smaller), PM10, NO2 (nitrogen dioxide), SO2 (sulfur dioxide), CO (carbon monoxide), O3 (ozone), and meteorological factors like temperature, humidity, wind speed, etc.

\subsection{Data Collection and Reliability}
The dataset is compiled from the official portal of the Government of India, the Central Pollution Control Board (CPCB). The CPCB is responsible for monitoring and controlling pollution, making the data collected from this source credible and trustworthy. The data collection process utilized Selenium, a web automation tool, for extracting and processing data from the CPCB website. This approach ensures the accuracy and reliability of the data\cite{CPCB:2023}.

\section{Dataset Details}

\subsection{Structure}
The dataset is structured with rows representing different cities and columns representing various air quality parameters and meteorological factors. and each city has different csv files in sub folders of states. 

\subsection{Size}
The dataset contains air quality data for 453 cities across India. It spans a time period from 2010 to 2023, covering 13 years of data. and 1.5GB of CSV files.

\subsection{Format}
The data is likely organized in a tabular format, with rows representing individual cities and columns representing different parameters such as PM2.5, PM10, NO2, SO2, CO, O3, temperature, humidity, wind speed, etc\cite{CPCB:2023}.

\subsection{Anomalies}
There are anomalies or missing values present in the dataset due to various reasons such as equipment malfunction, data collection errors, or incomplete records. These anomalies could affect the integrity of the dataset and might need to be addressed during data analysis.

\subsection{Origin}
The data has been sourced from the Central Control Room for Air Quality Management. It is collected and maintained by the Central Pollution Control Board (CPCB), which is an official agency responsible for monitoring and controlling pollution in India. The data collection process involved the use of Selenium, a web automation tool, to extract and process data from the CPCB website, ensuring the accuracy and reliability of the dataset\cite{CPCB:2023}.
